\documentclass{article}
\usepackage[utf8]{inputenc}
\usepackage[T1]{fontenc}
\usepackage[a4paper,total={6in,8in}]{geometry}
\usepackage{amsmath,amssymb,amsfonts}
\usepackage{graphicx}
\usepackage{hyperref}
\usepackage{enumitem}
\usepackage{listings}
\usepackage{xcolor}
\DeclareMathOperator{\Var}{Var}
\DeclareMathOperator{\Cov}{Cov}
\providecommand{\tightlist}{\setlength{\itemsep}{0pt}\setlength{\parskip}{0pt}}
\lstset{basicstyle=\ttfamily\scriptsize,breaklines=true,columns=fullflexible,
  keepspaces=true,showstringspaces=false,frame=single}
\hypersetup{colorlinks=true,urlcolor=blue,linkcolor=blue}
\setlength{\emergencystretch}{2em}

\title{Math 126 Problem 11: Random Sign Matrix}
\author{} % Put your name here.
\date{Spring 2027}
\begin{document}
\maketitle

\hypertarget{submission-instructions}{%
\subsection*{Submission Instructions}\label{submission-instructions}}

Please upload your work to Gradescope by \textbf{11:59pm on Monday, March 1, 2027 (Pacific time).}

\begin{itemize}
\item
  You are encouraged to discuss ideas and work with your classmates. However, you \emph{must acknowledge} your collaborators at the top of each solution on which you collaborated with others and you \emph{must write} your solutions independently.
\item
  Your solutions to theory questions must be written legibly, or typeset in LaTeX or markdown.
\item
  I recommend that you write your solutions in a Python notebook using Google Colab.
\item
  You should submit your solutions as a \emph{single PDF} via the assignment on Gradescope.
\end{itemize}

\textbf{Grading:} The point of the problem set is for \emph{you} to learn. To this end, I hope to disincentivize the use of LLMs by \emph{not} grading your work for correctness. Instead, you will grade your own work by comparing it to my solutions. This self-grade is due the Friday \emph{after} the problem set is due, also on Gradescope.

\noindent Original notebook: \href{https://colab.research.google.com/drive/1aW6LjXEWIVPulxRxWJs0Q1rrR6k4pq4W?usp=sharing}{Google Colab}.

% Source notebook cell 2
\hypertarget{random-sign-matrix}{%
\subsection*{Random Sign Matrix}\label{random-sign-matrix}}

Often practitioners prefer JL matrices with discrete random entries instead of Gaussians because they take less space to store and are easier to generate. We analyze one construction below.

Suppose that \(\Pi\) is a ``sign Johnson-Lindenstrauss matrix'' with \(n\) columns, \(k\) rows, and i.i.d. \(\pm 1\) entries scaled by \(1/\sqrt{k}\). In other words, each entry in the matrix has values \(-1/\sqrt{k}\) with probability \(1/2\) and \(1/\sqrt{k}\) with probability \(1/2\).

% Source notebook cell 3
\hypertarget{part-a}{%
\subsection*{Part A}\label{part-a}}

Prove that for any vector \(x \in \mathbb{R}^n\), \(E[\|\Pi x\|_2^2] = \|x\|_2^2\) and that \(\text{Var}[\|\Pi x\|_2^2] \leq \frac{2}{k} \|x\|_2^4\). This is the meat of the problem and will take some effort.

% Source notebook cell 4
\hypertarget{part-b}{%
\subsection*{Part B}\label{part-b}}

Use the above to prove that: \begin{align*}
\Pr \left( \left| \|\Pi x\|_2^2 - \|x\|_2^2 \right| \geq \epsilon \|x\|_2^2 \right) \leq \delta
\end{align*} as long as we choose \(k = O\left( \frac{1}{\delta \epsilon^2} \right)\).

\emph{Note: This bound almost matches the distributed JL lemma proven in class, but with a worse failure probability dependence of \(1/\delta\) in place of \(\log(1/\delta)\). With more work, it's possible to improve the dependence to \(\log(1/\delta)\) for the sign-JL matrix, but we won't do so here.}

\end{document}
